\section{Complete one-sided orbit optimization, including all hard cases}
\label{sec:complete-trust-region}

The generic optimization method in this section is classical trust-region
theory; see Mor\'e--Sorensen (1983) and Fortin--Wolkowicz (2004).  The extension
here is its exact application to the complete physical one-sided $F_3$ orbit,
with singular filters retained, together with a constructive formal proof.
In particular, a scalar minimization is not left as the answer: an explicitly
tested spectral endpoint or one uniquely admissible root of a degree-six
polynomial determines the result.

\subsection{Spectral algorithm for an arbitrary real quadratic}

Let $K=K^T\in\mathbb R^{3\times3}$, $k\in\mathbb R^3$, and $\tau\in\mathbb R$.
We seek the exact maximum of
\begin{equation}
 f(x)=\tau+2k^Tx-x^TKx,\qquad \|x\|\leq1.
 \label{eq:full-orbit-quadratic}
\end{equation}
Choose an orthogonal diagonalization
\[
 K=U\operatorname{diag}(d_1,d_2,d_3)U^T,\qquad h=U^Tk,
 \qquad d_{\min}=\min_i d_i,
 \qquad L=\max\{0,-d_{\min}\}.
\]
For $\ell\geq L$, define the partial reciprocal vector and its squared norm by
\[
 q_i(\ell)=
 \begin{cases}h_i/(d_i+\ell),&d_i+\ell\ne0,\\0,&d_i+\ell=0,\end{cases}
 \qquad \psi(\ell)=\sum_i q_i(\ell)^2.
\]
The zero-denominator convention alone does not make the endpoint admissible:
the range condition below must also hold.

\paragraph{Exact finite algorithm.}
First test the endpoint condition
\begin{equation}
 h_i=0\quad\hbox{whenever }d_i+L=0,
 \qquad \psi(L)\leq1.
 \label{eq:floor-admissibility}
\end{equation}
If it holds, set $\ell_*=L$.  If it fails, let $\ell_*>L$ be the unique solution of
\begin{equation}
 \sum_{i=1}^3\frac{h_i^2}{(d_i+\ell_*)^2}=1.
 \label{eq:regular-secular}
\end{equation}
This regular equation is equivalently the polynomial equation
\begin{align}
 P(\ell)={}&h_1^2(d_2+\ell)^2(d_3+\ell)^2
 +h_2^2(d_1+\ell)^2(d_3+\ell)^2\notag\\
 &+h_3^2(d_1+\ell)^2(d_2+\ell)^2
 -(d_1+\ell)^2(d_2+\ell)^2(d_3+\ell)^2=0.
 \label{eq:secular-polynomial}
\end{align}
Its degree is exactly six, with leading coefficient $-1$, including when
$h=0$ or eigenvalues coincide.  Consequently the endpoint and the real roots
of $P$ form a set of at most seven candidates.  The endpoint test
\eqref{eq:floor-admissibility} and the regular restriction $\ell>L$ select
exactly one admissible candidate.  Clearing denominators introduces no
spurious candidate on that regular interval, because every denominator is
strictly positive there.  In the original coordinates the same regular
polynomial equation can be written
\[
 \|\operatorname{adj}(K+\ell I)k\|^2-\det(K+\ell I)^2=0.
\]
The spectral version is particularly convenient for explicitly resolving the
singular endpoint.

\paragraph{Exact maximum and every maximizing point.}
The answer is
\begin{equation}
 M=\tau+\ell_*+
 \sum_{i:\,d_i+\ell_*\ne0}\frac{h_i^2}{d_i+\ell_*}.
 \label{eq:complete-orbit-value}
\end{equation}
Put $q=q(\ell_*)$ and
$E_* = \{z\in\mathbb R^3:z_i=0\text{ whenever }d_i+\ell_*\ne0\}$.
The complete maximizing set is
\begin{equation}
 \left\{U(q+z):z\in E_*,\quad
 \begin{cases}
  \|z\|^2\leq1-\psi(\ell_*),&\ell_*=0,\\
  \|z\|^2=1-\psi(\ell_*),&\ell_*>0.
 \end{cases}\right\}.
 \label{eq:complete-orbit-maximizers}
\end{equation}
The endpoint range condition guarantees $h_i=0$ on $E_*$, and $q$ is
orthogonal to $E_*$.  Thus this description includes every degenerate case:

\begin{itemize}
\item If $\ell_*>L$, all shifted eigenvalues are positive, and the unique
maximizer is $Uq$, on the unit sphere.
\item If $\ell_*=0$, the original $K$ is positive semidefinite.  The
maximizers form the explicitly displayed affine ball in $\ker K$.
For $K\succ0$ this is the single unconstrained vertex $K^{-1}k$.
\item If $\ell_*=L>0$, then $d_{\min}<0$ and $h$ is orthogonal to the
lowest-eigenvalue eigenspace.  This is the hard case.  The missing radius
is supplied by a vector in that eigenspace, giving precisely the sphere
in \eqref{eq:complete-orbit-maximizers}.  When its radius is zero the
maximizer is nevertheless unique.
\item If $k=0$, the formula reduces to $M=\tau+\max\{0,-d_{\min}\}$.
For $d_{\min}>0$ the maximizer is zero; for $d_{\min}=0$ it is the unit
ball in the kernel; for $d_{\min}<0$ it is the unit sphere in the
lowest-eigenvalue eigenspace.  This includes $K=0$.
\end{itemize}

\subsection{Constructive proof, without assumed optimality conditions}

In the eigenbasis, write $y=U^Tx$ and
$f_D(y)=\tau+\sum_i(2h_i y_i-d_i y_i^2)$.
We construct $p$ and $\ell$ satisfying
\begin{equation}
 \ell\geq0,\quad d_i+\ell\geq0,\quad
 (d_i+\ell)p_i=h_i,\quad
 \|p\|\leq1,\quad \ell(1-\|p\|^2)=0.
 \label{eq:constructed-certificate}
\end{equation}
These conditions are conclusions of the construction, not hypotheses
substituted for the optimization problem.

Suppose first that \eqref{eq:floor-admissibility} holds.
At $\ell=L=0$, take $p=q(L)$.  At $\ell=L>0$, choose an index $j$ for
which $d_j+L=0$, and set
\[
 p=q(L)+\sqrt{1-\psi(L)}\,e_j.
\]
The range condition gives $q_j(L)=h_j=0$, so $\|p\|=1$ and every equation
in \eqref{eq:constructed-certificate} holds.  This directly treats the
hard case, including repeated lowest eigenvalues and a zero linear term.

If the endpoint is inadmissible, a regular root exists by a finite bracket.
When the range condition holds but $\psi(L)>1$, choose $a=L$.
If instead $d_j+L=0$ and $h_j\ne0$ for some $j$, choose
$a=L+|h_j|/2$.  The $j$th summand of $\psi(a)$ is then exactly $4$,
so $\psi(a)>1$.  In either case all denominators at $a$ are nonnegative,
and any zero denominator has zero numerator.  The function $\psi$ is
therefore continuous on $[a,\infty)$.  For the explicit endpoint
\[
 b=a+1+\|h\|
\]
we have
\[
 \psi(b)\leq\frac{\|h\|^2}{(1+\|h\|)^2}<1.
\]
The intermediate value theorem gives a root $\ell\in(a,b)$, hence
$\ell>L$.  Now take $p=q(\ell)$; this proves
\eqref{eq:constructed-certificate} in the remaining case.  No compactness
or differentiability assumption about an unknown optimizer was needed.

For every such constructed certificate and every $y\in\mathbb R^3$,
direct expansion yields the exact gap identity
\begin{equation}
 f_D(p)-f_D(y)=
 \sum_i(d_i+\ell)(y_i-p_i)^2+\ell(1-\|y\|^2).
 \label{eq:trust-exact-gap}
\end{equation}
Both terms are nonnegative on the ball, proving global optimality.
Equality holds exactly when $y_i=p_i$ for every strictly positive
$d_i+\ell$, together with $\ell(1-\|y\|^2)=0$.
This proves \eqref{eq:complete-orbit-maximizers}.  Multiplying each
stationary equation by $p_i$ and using complementarity gives
\eqref{eq:complete-orbit-value}.

For completeness, the admissible shift is also unique in the hard and flat
cases.  Let $(p,\ell)$ and $(q,m)$ be two certificates.  Each point is
globally maximizing, so equality in \eqref{eq:trust-exact-gap} gives
$(d_i+\ell)q_i=h_i$ and $\ell(1-\|q\|^2)=0$.
Subtracting $(d_i+m)q_i=h_i$ yields $(m-\ell)q_i=0$ for every $i$.
If $m\ne\ell$, then $q=0$, and the two complementarity equations force
$m=\ell=0$, a contradiction.  Thus the algorithm selects one shift even
when it selects a continuum of maximizing vectors.

Equivalently, in the original coordinates, a point $x$ is globally
maximizing if and only if it belongs to the unit ball and there exists
$\ell\geq0$ with
\[
 K+\ell I\succeq0,\qquad (K+\ell I)x=k,
 \qquad \ell(1-\|x\|^2)=0.
\]
The proof above establishes existence of this global certificate for every
real symmetric $K$, rather than assuming this characterization.

\subsection{Exact application to all canonical quantum-filter orbits}

Let $\sigma$ be any actual two-qubit density matrix with Bob marginal $I/2$,
Alice Bloch vector $c\in\mathbb R^3$, and correlation matrix $S$.
Set
\begin{equation}
 \tau=\|S\|_F^2,\qquad k=S^Tc,\qquad
 K=(\tau-\|c\|^2)I-S^TS.
 \label{eq:physical-trust-data}
\end{equation}
For every complex $2\times2$ filter $F$ normalized by
$\operatorname{Tr}(F^\dagger F)=2$, there is an $x$ in the closed unit ball
such that $F^\dagger F=I+x\cdot\sigma$.  The actual normalized output
$\rho_F=(I\otimes F)\sigma(I\otimes F^\dagger)$ satisfies the exact identity
\begin{equation}
 F_3(\rho_F)^2=\tau+2k^Tx-x^TKx.
 \label{eq:physical-complete-objective}
\end{equation}
Conversely every such $x$ is realized by the positive filter
$F=(I+x\cdot\sigma)^{1/2}$.  These statements hold at every joint-state
rank and include singular filters.

Apply the finite algorithm to \eqref{eq:physical-trust-data}.
The exact maximum over all normalized complex one-sided filters is the
number $M$ in \eqref{eq:complete-orbit-value}, and the exact maximum of
$F_3$ is $\sqrt M$.  Every maximizing effect, and no other effect, is
\[
 F^\dagger F=I+\bigl(U(q+z)\bigr)\cdot\sigma,
\]
with $z$ exactly as in \eqref{eq:complete-orbit-maximizers}.  All complex
filters with a given maximizing effect are included; no Hermiticity
restriction is imposed on $F$.  A positive square-root filter explicitly
realizes each maximizing effect.

This completes the orbit optimization even when its maximum is at or below
the steering threshold.  The earlier violation criterion is recovered as
the special case in which the physical sphere bound $F_3^2\leq1$ forces
an interior optimizer with $K\succ0$.

\subsection{Invertible filters: supremum versus attainment}

For normalized filters,
\[
 \det(F^\dagger F)=1-\|x\|^2,
\]
so invertibility is equivalent to $\|x\|<1$.  The exact supremum over
invertible filters is still $M$: if $x_*$ is any closed-ball maximizer,
the effects $t x_*$ with $0\leq t<1$ are realized by invertible positive
filters, and their values tend to $M$ as $t\uparrow1$.

An invertible filter attains that supremum if and only if
\begin{equation}
 \ell_*=0\quad\hbox{and}\quad\psi(\ell_*)<1.
 \label{eq:invertible-attainment}
\end{equation}
Indeed positive shift forces every optimizer onto the unit sphere.
At zero shift the squared norm of $q+z$ is
$\psi(0)+\|z\|^2$, and an interior point exists exactly under
\eqref{eq:invertible-attainment}.  In particular, an endpoint maximum
can be attained by singular filters while only being approached by
invertible ones.  This distinction is not removed by assuming that the
input joint state has full rank.

\subsection{State rank, quadratic rank, and filter rank are different}
Two full-rank physical examples show why the singular cases cannot be
removed by imposing positivity on the input state.

First take
\[
 \sigma=\frac{I+cZ}{2}\otimes\frac I2,\qquad 0<c<1.
\]
Here $S=0$, $k=0$, and $K=-c^2I$. The objective is
$f(x)=c^2\|x\|^2$. Its maximum is $c^2$, every unit vector maximizes,
and the spectral shift is $\ell_*=c^2$. This is a hard case with a
three-dimensional nullspace of $K+\ell_*I$. Every maximizing filter is
singular, and the same value is only a supremum for invertible filters,
even though the input density matrix has rank four.

Next take
\[
 \sigma=\frac14(I+tX\otimes X),\qquad 0<t<1.
\]
This state also has rank four, but now $k=0$ and
$K=\operatorname{diag}(0,t^2,t^2)$. The exact maximum is $t^2$ and every
point $x=(s,0,0)$ with $|s|\leq1$ is a maximizer. Thus the complete
maximizing effect set contains both invertible and singular filters.
These are direct specializations of the universal checked theorem.
The exact rational certificate file includes both types of example.

\paragraph{Verification scope.}
The accompanying formal-coverage manifest identifies the exact Lean
declarations and compiler evidence for the results above.  Numerical
comparisons are reproducibility checks only; no numerical result is a
logical premise of a formal theorem.
